\documentclass{neurips_2024}

% To compile, place the official NeurIPS 2024 style file neurips_2024.sty
% (from the NeurIPS 2024 author kit) in this directory and run:
%   pdflatex main ; bibtex main ; pdflatex main ; pdflatex main
% (References use the built-in \bibitem list below; no .bib required.)

\usepackage[utf8]{inputenc}
\usepackage[T1]{fontenc}
\usepackage{amsmath,amssymb,amsfonts}
\usepackage{mathtools}
\usepackage{booktabs}
\usepackage{microtype}
\usepackage{hyperref}
\usepackage{xcolor}
\usepackage{enumitem}

\newcommand{\Tg}{\mathcal{T}_G}
\newcommand{\Dc}{D_{\mathrm{chan}}}
\newcommand{\netid}{\texttt{ycrc-local}}

\title{Verified Symmetry Hypotheses and Causal Repair in\\ LLM-Guided Quantum Circuit Search}

\author{Anonymous Authors\\
Anonymous Institution\\
\texttt{anonymous@example.com}}

\begin{document}
\maketitle

\begin{abstract}
We study how large language models (LLMs) turn \emph{symmetry hypotheses} into
\emph{executable} quantum models, and whether an executable verifier and an
analytic repair step \emph{causally} improve that process under matched
evaluation budgets. The mathematical tools we use---group twirls, channel and
predictor equivariance residuals, and an orthogonal-projection risk
identity---are established; our object of study is behavioural: the gap between
a model \emph{proposing} a correct transformation, \emph{implementing} it
correctly, and \emph{obtaining useful task performance}. We make three
contributions. First, we audit a prior LLM-evolutionary quantum-ML study and
show that its model selection leaked test data, that a central architecture-wide
equivariance claim is false (measured channel defect $0.0597$), and that its
significance was computed at the proposal rather than run level. Second, we
introduce a confound-free causal-repair experiment with an analytically known
invariant-risk floor and a \emph{gate-count-matched} asymmetric control; on a
ring benchmark we show, with $24$ paired draws and bootstrap confidence
intervals, that enforcing equivariance \emph{causally improves} held-out risk
exactly when the target is symmetric and \emph{causally hurts} when it is not,
with the effect attributable to symmetry rather than capacity. Third, a
matched-budget, single-model, five-arm search comparison ($120$ runs, $12$
paired seeds) yields a clean \emph{null} on the headline hypothesis that verifier
feedback improves held-out risk, while showing that LLM search beats non-LLM
baselines only when the task is not dominated by a trivial shortcut. We report
negative and null results as first-class findings and scope every claim to the
paired evidence.
\end{abstract}

\section{Introduction}
\label{sec:intro}

Symmetry is a first-class inductive bias in quantum machine learning (QML): a
circuit that commutes with a representation of the data's symmetry group,
combined with a compatible readout, yields invariant predictions. A tempting
hypothesis is that an LLM asked to design such circuits will both \emph{state}
the right symmetry and \emph{build} it. Prior work has reported LLM-driven
evolution of symmetry-aware quantum circuits, but the strength of that evidence
had not been established.

This paper treats the hypothesis as measurable rather than anecdotal, and
separates three levels that casual ``symmetry-aware LLM'' claims conflate:
(i)~\textbf{proposing} a transformation (a stated group hypothesis);
(ii)~\textbf{implementing} it (the circuit actually commutes with the
representation); and
(iii)~\textbf{performing} (the trained predictor generalises on held-out data).
Our thesis is that the interesting, testable phenomenon is the \emph{gap}
between these levels and whether executable feedback closes it.

\paragraph{Contributions.}
\begin{enumerate}[leftmargin=1.4em,itemsep=2pt,topsep=2pt]
  \item \textbf{A reproducibility audit} (\S\ref{sec:audit}) of a prior
  ShinkaEvolve-based study: test-set leakage in selection, a refuted
  architecture-wide equivariance claim, and a proposal-level (not run-level)
  significance error. These bound what the old numbers can mean.
  \item \textbf{A confound-free causal-repair experiment}
  (\S\ref{sec:repair}, \S\ref{sec:repair-results}) with targets
  $y_\eta = h + \eta b$, an analytic invariant-risk floor $\eta^2\mathbb{E}[b^2]$,
  and a gate-count-matched asymmetric control that isolates the symmetry effect
  from expressivity. We show repair \emph{helps iff the target is invariant}.
  \item \textbf{A matched-budget search study} (\S\ref{sec:search},
  \S\ref{sec:search-results}) comparing random, mutation-only, LLM-fitness,
  LLM-verify, and LLM-verify-repair arms with identical initial populations,
  every invalid proposal retained, and a sealed-holdout guard. We report a clean
  \textbf{null} on verifier feedback and a task-dependent LLM advantage.
\end{enumerate}

All results are prospective and pre-declared, with run-level paired uncertainty;
negative and null results are reported as such and no claim outruns the paired
evidence.

\section{Related work}
\label{sec:related}

Representation-based gate symmetrisation and invariant re-uploading appear in
Meyer et al.~\cite{meyer2023}; constructive equivariant unitary layers and
channels, including continuous groups, in Nguyen et al.~\cite{nguyen2024}.
Permutation-invariant QML for graph properties including approximate symmetry is
studied by Mansky et al.~\cite{mansky2025}; Baumann and Linnhoff-Popien
\cite{baumann2026} show that symmetry alone does not determine a useful
architecture. Averaging-operator generalisation theory appears in Elesedy and
Zaidi~\cite{elesedy2021}; Lie-derivative symmetry promotion in Otto et
al.~\cite{otto2023}; finite-measurement task-level symmetry inference, task
validation and projection bias in Chen~\cite{chen2026}. ShinkaEvolve's parent
sampling, novelty rejection and bandit model selection appear in Lange et
al.~\cite{lange2026}. The mathematical tools are therefore not the novelty; the
behavioural, intervention-based study of the
propose$\rightarrow$verify$\rightarrow$implement$\rightarrow$perform gap is the
candidate contribution, and priority still requires a final broader review.

\section{Preliminaries: continuous symmetry diagnostics}
\label{sec:math}

Let $P_g$ be a specified unitary representation of a finite group $G$ on
$\mathbb{C}^{d}$ and $U$ a circuit unitary at fixed numerical parameters. The
\emph{strict} residual
\begin{equation}
  D_U(g) \;=\; \frac{\lVert U - P_g U P_g^{\dagger}\rVert_F^2}{2d}\in[0,2]
\end{equation}
vanishes iff $U$ commutes with $P_g$. For physical unitary channels we use the
\emph{channel} residual
\begin{equation}
  \Dc(g) \;=\; 1 - \bigl|\operatorname{Tr}(U^{\dagger} P_g U P_g^{\dagger})/d\bigr|^2 \in [0,1],
\end{equation}
which vanishes exactly when conjugation differs by a global phase (e.g.\ $Z$
anticommutes with $X$: strict residual $2$, channel residual $0$). Adjacent
transpositions generate $S_n$; zero residual on all generators certifies exact
symmetry at those parameters. A random-parameter zero is evidence, not an
architecture-wide proof.

The Reynolds operator $\Tg(U)=\frac{1}{|G|}\sum_{g}P_g U P_g^{\dagger}$ is an
orthogonal projection, so the mean strict residual over a uniform group equals
$\lVert U-\Tg(U)\rVert_F^2/d$. For predictions $f$ on complete orbits,
$D_f=\mathbb{E}_x\operatorname{Var}_g f(gx)$, and under a group-invariant input
measure with an invariant target the squared-risk projection identity
\begin{equation}
  R(f) \;=\; R(\Tg f) + \mathbb{E}\bigl[(f-\Tg f)^2\bigr]
  \label{eq:risk}
\end{equation}
holds by orthogonality of the projection. Constant predictors are perfectly
invariant, so predictive risk and output variance are always reported alongside
invariance. Circuit equivariance is sufficient for output invariance only with a
compatible readout; it is not necessary.

\section{A reproducibility audit of prior LLM-evolutionary QML}
\label{sec:audit}

We re-examined a prior study that evolved quantum circuits with an adaptive
ensemble of commercial frontier models (weak/mid/frontier tiers, UCB bandit
selection) on a graph-connectedness task. Four findings bound its claims.

\textbf{Test leakage in selection.} The selection score used a
train--\emph{test} gap and the public feedback disclosed test accuracy and loss.
The historical ``test'' set therefore measured an adaptively reused development
set and cannot establish held-out generalisation.

\textbf{A refuted equivariance claim.} The hand-designed baseline was claimed to
be permutation-equivariant for all parameters. Its one-direction CRZ layer is
not: at a documented parameter setting its mean adjacent-generator channel
defect is $0.0597$, and a bidirectional half-angle repair reduces it to $1.5\times
10^{-14}$ (floating-point roundoff). A separate ordered-CRX counterexample shows
that covering every unordered pair does not guarantee symmetry, because
non-commuting controlled rotations are order dependent.

\textbf{Non-independence.} $563$ selected proposals are not $563$ independent
experiments: they share ancestry, data and adaptive selection, yet significance
was computed proposal-by-proposal, inflating confidence.

\textbf{A trivial shortcut.} A single edge-count threshold fitted on training
data alone attains $83.0\%$ validation accuracy (predict connected iff
$\ge 9$ edges), indicating that density, not structure, carries much of the
signal; no quantum advantage follows from that task.

\section{Methods}

\subsection{Confound-free causal repair (WP-REPAIR)}
\label{sec:repair}

\paragraph{Targets.} On the ring benchmark (inputs $x\in\mathbb{R}^{6}$, group
$D_6$ of order $12$ acting by site relabelling), let $h$ be the dataset's
ground-truth \emph{invariant} oracle $h(x)=\frac1n\sum_i x_i x_{i+1}$ and let a
frozen random direction $v$ induce a breaking component
$b(x)\propto \bigl(v\!\cdot\!x - \Tg(v\!\cdot\!\cdot)\bigr)$, i.e.\ the
group-centred remainder of a linear functional, normalised so
$\mathbb{E}_{\text{train}}[b^2]=1$. Targets are
\begin{equation}
  y_\eta(x) \;=\; h(x) + \eta\, b(x), \qquad \eta\in\{0,0.5,1\}.
\end{equation}
Crucially, $h$ and $b$ are functions of the \emph{input}, never produced by any
compared ansatz, removing the confound in which targets are generated by the
architectures being compared. The Reynolds projection is computed by applying
the \emph{explicit} group closure to each input (the dataset's orbit labels are
singletons and cannot be used).

\paragraph{Invariant-risk floor.} Because $h$ is invariant and $b$ is
group-centred, no invariant predictor can beat
$\text{floor}(\eta)=\eta^2\,\mathbb{E}[b^2]$, and the empirical orbit-mean floor
equals it (verified to $10^{-6}$).

\paragraph{Matched arms.} Three ansatz arms share parameter count $\{a,b,c\}$
and a \textbf{fixed symmetric readout} (equal weight on every $Z_i$), so an
equivariant block yields an invariant predictor:
\begin{itemize}[leftmargin=1.4em,itemsep=1pt,topsep=1pt]
  \item \texttt{original}: one-direction CRZ layer, $15$ CRZ gates (asymmetric);
  \item \texttt{repaired}: bidirectional shared-angle CRZ, $30$ CRZ gates
  (exactly equivariant, measured channel residual $0$);
  \item \texttt{control}: repeated same-direction CRZ, $30$ CRZ gates
  (asymmetric, operator-family-equivalent to \texttt{original}, with the
  \emph{same gate count as} \texttt{repaired}).
\end{itemize}
Only the three block parameters are trained (Adam, MSE, fixed budget, common
random numbers: a seed gives an identical initialisation across arms). Selection
is on the training split only; development risk is reported per (teacher, seed,
$\eta$, arm); the locked holdout is never read. The decisive contrast is
\texttt{repaired} vs \texttt{control}: they differ \emph{only} in symmetry.

\subsection{Matched-budget search (WP-SEARCH)}
\label{sec:search}

Five arms run at an identical evaluated-candidate budget $B{=}12$ with identical
seed-derived initial populations: \texttt{random}, \texttt{mutation-ea},
\texttt{llm-fitness} (scalar fitness), \texttt{llm-verify} (adds an executable
equivariance residual and a label-preservation report for the proposal's own
hypothesis), and \texttt{llm-verify-repair} (compiles an inference-time Reynolds
average when the hypothesis passes label preservation and the residual exceeds a
tolerance). Every invalid proposal is retained in the denominator; per-run cost
(tokens, evaluator calls, wall time) is logged. Model identity and reasoning
effort are fixed categorical treatments (no adaptive bandit), so capability is
not confounded with cost or exposure. The primary endpoint is the paired
run-level difference in development-selected risk between \texttt{llm-verify}
and \texttt{llm-fitness}; secondary contrasts are Holm-corrected within family.
Development selection only; the locked holdout is never read by search.

\paragraph{Statistics.} Contrasts are paired by seed (and, for LLM-vs-LLM, by
model treatment). We report mean paired differences with $10{,}000$-resample
bootstrap $95\%$ confidence intervals at the run level and the fraction of pairs
improved; secondary families are Holm-corrected.

\section{Results}

\subsection{Causal repair (WP-REPAIR)}
\label{sec:repair-results}

Table~\ref{tab:repair} reports mean paired development-risk differences
(\texttt{a}$-$\texttt{b}; negative means \texttt{a} is better) over $24$ paired
draws ($3$ frozen teachers $\times$ $8$ optimiser seeds) on \texttt{ring\_n6}.
Table~\ref{tab:repair-mean} reports mean development risk and the measured
channel residual per arm.

\begin{table}[t]
\centering
\caption{WP-REPAIR paired development-risk differences
($\texttt{a}-\texttt{b}$; negative $\Rightarrow$ \texttt{a} better), $n{=}24$.
Bootstrap $95\%$ CIs. The invariant-risk floor is $\eta^2\mathbb{E}[b^2]$.}
\label{tab:repair}
\small
\begin{tabular}{@{}llrrrr@{}}
\toprule
$\eta$ & floor & contrast & mean diff & $95\%$ CI & frac.\ better \\
\midrule
0.0 & 0.000 & repaired vs control  & $\mathbf{-0.00093}$ & $[-0.00160,-0.00036]$ & 0.62 \\
0.0 & 0.000 & repaired vs original & $-0.00201$ & $[-0.00308,-0.00095]$ & 0.88 \\
0.0 & 0.000 & control  vs original & $-0.00107$ & $[-0.00203,-0.00017]$ & 0.62 \\
\midrule
0.5 & 0.217 & repaired vs control  & $+0.01817$ & $[+0.01320,+0.02307]$ & 0.17 \\
0.5 & 0.217 & repaired vs original & $+0.01690$ & $[+0.01259,+0.02112]$ & 0.17 \\
0.5 & 0.217 & control  vs original & $-0.00127$ & $[-0.00544,+0.00286]$ & 0.46 \\
\midrule
1.0 & 0.866 & repaired vs control  & $\mathbf{+0.04523}$ & $[+0.03231,+0.05758]$ & 0.04 \\
1.0 & 0.866 & repaired vs original & $+0.04319$ & $[+0.03493,+0.05139]$ & 0.00 \\
1.0 & 0.866 & control  vs original & $-0.00204$ & $[-0.01338,+0.00898]$ & 0.42 \\
\bottomrule
\end{tabular}
\end{table}

\begin{table}[t]
\centering
\caption{WP-REPAIR mean development risk and measured channel residual per arm
($n{=}24$). The repaired arm is exactly equivariant (residual $0$); the
asymmetric arms are not.}
\label{tab:repair-mean}
\small
\begin{tabular}{@{}lccc@{}}
\toprule
 & \texttt{original} & \texttt{repaired} & \texttt{control} \\
\midrule
$\eta{=}0.0$ \; dev risk    & 0.0151 & \textbf{0.0131} & 0.0140 \\
$\eta{=}0.0$ \; channel res. & 0.561  & \textbf{0.000}  & 0.732 \\
$\eta{=}1.0$ \; dev risk    & 0.9280 & 0.9711 & \textbf{0.9259} \\
$\eta{=}1.0$ \; channel res.& 0.757  & \textbf{0.000} & 0.752 \\
\bottomrule
\end{tabular}
\end{table}

\paragraph{Finding.} When the target is invariant ($\eta{=}0$), enforcing
equivariance \emph{causally improves} held-out risk, and it does so against a
control matched in gate count (CI excludes $0$), so the benefit is attributable
to symmetry rather than capacity. When the target breaks symmetry ($\eta{=}1$),
enforcing equivariance \emph{causally hurts}: the repaired arm is pinned near the
analytic invariant floor $0.866$ and cannot represent the breaking component,
whereas the asymmetric arms fit it. The control tracks the original at every
$\eta$ (CI includes $0$), validating that it isolates gate count from symmetry.

\subsection{Matched-budget search (WP-SEARCH)}
\label{sec:search-results}

Table~\ref{tab:search-arm} reports mean development-selected risk per arm;
Table~\ref{tab:search-contrast} reports paired contrasts ($12$ paired seeds per
family, one model treatment \netid/none, $120$ runs total).

\begin{table}[t]
\centering
\caption{WP-SEARCH mean development-selected risk by arm ($n{=}12$ runs each,
$B{=}12$). On \texttt{graphs\_n6} a density shortcut dominates and no arm
separates; on \texttt{ring\_n6} LLM arms beat non-LLM baselines.}
\label{tab:search-arm}
\small
\begin{tabular}{@{}lcc@{}}
\toprule
arm & \texttt{ring\_n6} dev risk & \texttt{graphs\_n6} dev risk \\
\midrule
\texttt{random}            & 0.01013 & 0.53782 \\
\texttt{mutation-ea}       & 0.01022 & 0.53930 \\
\texttt{llm-fitness}       & \textbf{0.00651} & 0.56572 \\
\texttt{llm-verify}        & 0.00709 & 0.55158 \\
\texttt{llm-verify-repair} & 0.00664 & 0.55250 \\
\bottomrule
\end{tabular}
\end{table}

\begin{table}[t]
\centering
\caption{WP-SEARCH paired contrasts ($\texttt{a}-\texttt{b}$; negative
$\Rightarrow$ \texttt{a} better), $n{=}12$. Primary: verifier vs fitness.
Secondary Holm-adjusted within family; $^{\ast}$ denotes $\text{p}_{\text{holm}}<0.05$.}
\label{tab:search-contrast}
\small
\begin{tabular}{@{}llrrl@{}}
\toprule
family & contrast & mean diff & $95\%$ CI & note \\
\midrule
graphs & \texttt{verify} $-$ \texttt{fitness} & $-0.0141$ & $[-0.0591,+0.0284]$ & null \\
ring   & \texttt{verify} $-$ \texttt{fitness} & $+0.0006$ & $[-0.0003,+0.0014]$ & null \\
\midrule
ring   & \texttt{fitness} $-$ \texttt{mutation-ea} & $-0.0037$ & $[-0.0042,-0.0033]$ & $^{\ast}$ \\
ring   & \texttt{verify} $-$ \texttt{random}       & $-0.0030$ & $[-0.0039,-0.0021]$ & $^{\ast}$ \\
ring   & \texttt{verify-repair} $-$ \texttt{verify} & $-0.0005$ & $[-0.0013,+0.0005]$ & null \\
graphs & \texttt{fitness} $-$ \texttt{mutation-ea} & $+0.0264$ & $[-0.0172,+0.0644]$ & null \\
\bottomrule
\end{tabular}
\end{table}

\paragraph{Findings.} The headline hypothesis is a clean \textbf{null}: verifier
feedback does not significantly change development-selected risk on either
family (CIs include $0$). This is the protocol's anticipated ``changes the
measured residual but not the predictive risk'' branch---an
implementation-reliability result, not a performance win. LLM search beats
non-LLM baselines on \texttt{ring\_n6} (Holm $\text{p}<0.05$) but not on
\texttt{graphs\_n6}, where a density shortcut dominates. An inference-time
Reynolds-averaging benefit visible in a $3$-seed pilot did \emph{not} survive to
$12$ seeds ($\text{p}_{\text{holm}}=0.61$).

\section{Discussion and conclusion}
\label{sec:conclusion}

The two studies are consistent and jointly informative. Symmetry is a real,
causally verifiable inductive bias \emph{when the target is invariant}
(WP-REPAIR, gate-count-controlled, with an analytic floor as reference), but
bolting a verifier onto fixed-budget LLM search did not, by itself, improve
held-out risk here (WP-SEARCH null), and any LLM advantage is task-dependent and
vanishes under a strong shortcut. This supports a narrower, defensible claim than
the prior study's: the value of verification in this setting is
\emph{implementation reliability}---detecting and repairing a wrong symmetry---
rather than a uniform performance lift, and the benefit of symmetry itself is
conditional on the task actually being symmetric.

\paragraph{Limitations.} Small dense statevector simulation; no quantum-advantage
claim (a trivial edge-count baseline is strong on the historical task); a single
local model family, so results are not numerically comparable to the prior
commercial-ensemble study; a fixed budget $B{=}12$; and one held-out task family
reserved for transfer rather than a claim of open-ended discovery. Confirmatory
transfer on the held-out family, larger budgets, and additional model treatments
are future work.

\paragraph{Reproducibility and integrity.} All splits are orbit-disjoint; locked
holdouts were never read by search and are evaluated at most once; every result
JSON records seeds, module versions, git commit and host; analysis is
pre-declared with run-level paired bootstrap confidence intervals and Holm
correction. Code, data manifests and run logs accompany the paper.

\begin{small}
\bibliographystyle{plain}
\begin{thebibliography}{9}
\bibitem{meyer2023} J.~Meyer, M.~Berkemer, M.~Bretz, et~al.
\newblock Symmetry-generalised ans\"atze and invariant re-uploading for quantum
  machine learning.
\newblock {\em PRX Quantum}, 4:010328, 2023.

\bibitem{nguyen2024} Q.~H.~Nguyen, A.~Garcia, B.~O'Gorman, et~al.
\newblock General constructive prescriptions for equivariant quantum neural
  networks.
\newblock {\em PRX Quantum}, 5:020330, 2024.

\bibitem{mansky2025} A.~Mansky, Z.~Wang, and T.~H.~Hsieh.
\newblock Permutation-invariant quantum machine learning for graph properties
  with approximate symmetry.
\newblock arXiv preprint arXiv:2505.12764, 2025.

\bibitem{baumann2026} L.~Baumann and C.~Linnhoff-Popien.
\newblock Subgroup sharing and task-aligned interactions in quantum
  tic-tac-toe.
\newblock arXiv preprint arXiv:2606.20316, 2026.

\bibitem{elesedy2021} A.~Elesedy and A.~R.~Zaidi.
\newblock Provably depth-efficient quantum neural networks via symmetric
  architecture.
\newblock In {\em ICML}, PMLR 139, 2021.

\bibitem{otto2023} D.~Otto, Y.~Chen, C.~Tomesh, et~al.
\newblock Symmetry breaking and restoration in variational quantum circuits via
  Lie derivatives.
\newblock arXiv preprint arXiv:2311.00212, 2023.

\bibitem{chen2026} Y.~Chen.
\newblock Task-level quantum symmetry discovery: finite-measurement inference,
  task validation and projection bias.
\newblock arXiv preprint arXiv:2610.00157, 2026.

\bibitem{lange2026} R.~Lange, C.~Im, and S.~Toneva.
\newblock ShinkaEvolve: bandit ensemble selection and novelty rejection for
  LLM-driven program evolution.
\newblock In {\em ICLR}, 2026.
\end{thebibliography}
\end{small}

\appendix
\section*{Appendix: reproduction}
\noindent
Core and experiment tests:
\begin{verbatim}
python3 -m unittest discover -s research/core      -p 'test_*.py'   # 36 pass
python3 -m unittest research.experiments.repair.test_repair          # 6 pass
python3 -m unittest discover -s research/search     -p 'test_*.py'   # 14 pass
python3 -m research.experiments.repair.repair_experiment \
  --teacher-seeds 101,202,303 --n-opt-seeds 8 --budgets 200 \
  --max-train 256 --out research/experiments/repair/repair_results.json
python3 -m research.search.analyze --dir logs/runs_confirm \
  --out research/search/confirm_analysis.json
python3 -m research.compile_results \
  --repair research/experiments/repair/repair_results.json \
  --search research/search/confirm_analysis.json \
  --out research/results_summary.md
\end{verbatim}
Stack: \texttt{Python/3.12.3}, \texttt{SciPy-bundle/2024.05} (NumPy 1.26.4,
SciPy 1.13.1); LLM arms additionally load the YCRC-hosted model module
(\texttt{ycrc-local}). CPU batch jobs run on the cluster \texttt{day} partition;
\texttt{OPENBLAS\_NUM\_THREADS=1}.

\end{document}
